\begin{textsample}{POS Dim 1 – 2003 – Score 23.00 – 2003/t000259.txt}  \label{ex:f1_pos_003}
Last year, I told you \textbf{the} story, in seven minutes, of Project Orion, which was this very implausible technology that technically could have worked, but it had this one-year political window where it could have happened.
So it didn ' t happen.
It was a dream that did not happen.
This year I ' m going to tell you \textbf{the} story of \textbf{the} birth of digital computing.
This was a perfect introduction.
And it ' s a story that did work.
It did happen, and \textbf{the} machines are all around us.
And it was a technology that was inevitable.
If \textbf{the} people I ' m going to tell you \textbf{the} story about, if they hadn ' t done it, somebody else would have.
So, it was sort of \textbf{the} right idea at \textbf{the} right time.
This is Barricelli ' s \textbf{universe}.
This is \textbf{the} \textbf{universe} we live in now.
It ' s \textbf{the} \textbf{universe} in which these machines are now doing all these things, including changing \textbf{biology}.
I ' m starting \textbf{the} story with \textbf{the} first atomic bomb at Trinity, which was \textbf{the} Manhattan Project.
It was a little bit like TED : it brought a whole lot of very smart people together.
And three of \textbf{the} smartest people were Stan Ulam, Richard Feynman and John von Neumann.
And it was Von Neumann who said, after \textbf{the} bomb, he was working on something much more important than bombs : he ' s thinking about \textbf{computers}.
So, he wasn ' t only thinking about them ; he built one.
This is \textbf{the} machine he built.
He built this machine, and we had a beautiful demonstration of how this thing really works, with these little bits.
And it ' s an idea that goes way back. \textbf{The} first person to really explain that was Thomas Hobbes, who, in 1651, explained how arithmetic and logic are \textbf{the} same thing, and if you want to do \textbf{artificial} thinking and \textbf{artificial} logic, you can do it all with arithmetic.
He said you needed addition and subtraction.
Leibniz, who came a little bit later — this is 1679 — showed that you didn ' t even need subtraction.
You could do \textbf{the} whole thing with addition.
Here, we have all \textbf{the} binary arithmetic and logic that drove \textbf{the} \textbf{computer} revolution.
And Leibniz was \textbf{the} first person to really talk about building such a machine.
He talked about doing it with marbles, having gates and what we now call shift registers, where you shift \textbf{the} gates, drop \textbf{the} marbles down \textbf{the} tracks.
And that ' s what all these machines are doing, except, instead of doing it with marbles, they ' re doing it with electrons.
And then we jump to Von Neumann, 1945, when he sort of reinvents \textbf{the} whole same thing.
And 1945, after \textbf{the} war, \textbf{the} electronics existed to actually try and build such a machine.
So June 1945 — actually, \textbf{the} bomb hasn ' t even been dropped yet — and Von Neumann is putting together all \textbf{the} theory to actually build this thing, which also goes back to Turing, who, before that, gave \textbf{the} idea that you could do all this with a very brainless, little, finite state machine, just reading a tape in and reading a tape out. \textbf{The} other sort of genesis of what Von Neumann did was \textbf{the} difficulty of how you would predict \textbf{the} weather.
Lewis Richardson saw how you could do this with a cellular \textbf{array} of people, giving them each a little chunk, and putting it together.
Here, we have an \textbf{electrical} model illustrating a mind having a will, but capable of only two ideas.
And that ' s really \textbf{the} simplest \textbf{computer}.
It ' s basically why you need \textbf{the} qubit, because it only has two ideas.
And you put lots of those together, you get \textbf{the} essentials of \textbf{the} modern \textbf{computer} : \textbf{the} arithmetic unit, \textbf{the} central control, \textbf{the} memory, \textbf{the} recording medium, \textbf{the} input and \textbf{the} output.
But, there ' s one catch.
This is \textbf{the} fatal — you know, we saw it in starting these programs up. \textbf{The} instructions which govern this operation must be given in absolutely exhaustive detail.
So, \textbf{the} programming has to be perfect, or it won ' t work.
If you look at \textbf{the} origins of this, \textbf{the} classic history sort of takes it all back to \textbf{the} ENIAC here.
But actually, \textbf{the} machine I ' m going to tell you about, \textbf{the} Institute for Advanced Study machine, which is way up there, really should be down there.
So, I ' m trying to revise history, and give some of these guys more credit than they ' ve had.
Such a \textbf{computer} would open up \textbf{universes}, which are, at \textbf{the} present, outside \textbf{the} range of any instruments.
So it opens up a whole new world, and these people saw it. \textbf{The} guy who was supposed to build this machine was \textbf{the} guy in \textbf{the} middle, Vladimir Zworykin, from RCA.
RCA, in probably one of \textbf{the} lousiest business decisions of all time, decided not to go into \textbf{computers}.
But \textbf{the} first meetings, November 1945, were at RCA ' s offices.
RCA started this whole thing off, and said, you know, televisions are \textbf{the} future, not \textbf{computers}. \textbf{The} essentials were all there — all \textbf{the} things that make these machines run.
Von Neumann, and a logician, and a mathematician from \textbf{the} army put this together.
Then, they needed a place to build it.
When RCA said no, that ' s when they decided to build it in Princeton, where Freeman works at \textbf{the} Institute.
That ' s where I grew up as a kid.
That ' s me, that ' s my sister Esther, who ' s talked to you before, so we both go back to \textbf{the} birth of this thing.
That ' s Freeman, a long time ago, and that was me.
And this is Von Neumann and Morgenstern, who wrote \textbf{the}"Theory of Games."All these forces came together there, in Princeton.
Oppenheimer, who had built \textbf{the} bomb. \textbf{The} machine was actually used mainly for doing bomb calculations.
And Julian Bigelow, who took Zworkykin ' s place as \textbf{the} engineer, to actually figure out, using electronics, how you would build this thing. \textbf{The} whole gang of people who came to work on this, and women in front, who actually did most of \textbf{the} coding, were \textbf{the} first programmers.
These were \textbf{the} prototype \textbf{geeks}, \textbf{the} nerds.
They didn ' t fit in at \textbf{the} Institute.
This is a letter from \textbf{the} director, concerned about —"especially unfair on \textbf{the} matter of sugar."You can read \textbf{the} text.
This is hackers getting in trouble for \textbf{the} first time..
These were not theoretical physicists.
They were real soldering-gun type guys, and they actually built this thing.
And we take it for granted now, that each of these machines has billions of transistors, doing billions of cycles per second without failing.
They were using vacuum tubes, very narrow, sloppy \textbf{techniques} to get actually binary behavior out of these radio vacuum tubes.
They actually used 6J6, \textbf{the} common radio tube, because they found they were more reliable than \textbf{the} more expensive tubes.
And what they did at \textbf{the} Institute was publish every step of \textbf{the} way.
Reports were issued, so that this machine was \textbf{cloned} at 15 other places around \textbf{the} world.
And it really was.
It was \textbf{the} original microprocessor.
All \textbf{the} \textbf{computers} now are \textbf{copies} of that machine. \textbf{The} memory was in cathode ray tubes — a whole bunch of spots on \textbf{the} face of \textbf{the} tube — very, very sensitive to electromagnetic disturbances.
So, there ' s 40 of these tubes, like a V-40 engine running \textbf{the} memory. \textbf{The} input and \textbf{the} output was by teletype tape at first.
This is a wire drive, using bicycle wheels.
This is \textbf{the} archetype of \textbf{the} hard \textbf{disk} that ' s in your machine now.
Then they switched to a magnetic \textbf{drum}.
This is modifying IBM equipment, which is \textbf{the} origins of \textbf{the} whole data-processing industry, later at IBM.
And this is \textbf{the} beginning of \textbf{computer} \textbf{graphics}. \textbf{The}"Graph ' g-Beam Turn On."This next \textbf{slide}, that ' s \textbf{the} — as far as I know — \textbf{the} first digital bitmap display, 1954.
So, Von Neumann was already off in a theoretical cloud, doing \textbf{abstract} sorts of studies of how you could build reliable machines out of unreliable components.
Those guys \textbf{drinking} all \textbf{the} tea with sugar in it were writing in their logbooks, trying to get this thing to work, with all these 2, 600 vacuum tubes that failed half \textbf{the} time.
And that ' s what I ' ve been doing, this last six months, is going through \textbf{the} logs."Running time : two minutes.
Input, output : 90 minutes."This includes a large amount of human error.
So they are always trying to figure out, what ' s machine error?
What ' s human error?
What ' s code, what ' s hardware?
That ' s an engineer gazing at tube number 36, trying to figure out why \textbf{the} memory ' s not in focus.
He had to focus \textbf{the} memory — seems OK.
So, he had to focus each tube just to get \textbf{the} memory up and running, let alone having, you know, software problems."No use, went home.""Impossible to follow \textbf{the} damn thing, where ' s a directory?"So, already, they ' re complaining about \textbf{the} manuals :"before closing down in disgust...""\textbf{The} General Arithmetic : Operating Logs."Burning lots of midnight \textbf{oil}."MANIAC,"which became \textbf{the} acronym for \textbf{the} machine, Mathematical and Numerical Integrator and Calculator,"lost its memory.""MANIAC regained its memory, when \textbf{the} power went off.""Machine or human?""Aha!"So, they figured out it ' s a code problem."Found trouble in code, I hope.""Code error, machine not guilty.""Damn it, I can be just as stubborn as this thing.""And \textbf{the} dawn came."So they ran all night.
Twenty-four hours a day, this thing was running, mainly running bomb calculations."Everything up to this point is wasted time.""What ' s \textbf{the} use?
Good night.""Master control off. \textbf{The} hell with it.
Way off.""Something ' s wrong with \textbf{the} air conditioner — smell of burning V-belts in \textbf{the} air.""A short — do not turn \textbf{the} machine on.""IBM machine putting a tar-like substance on \textbf{the} cards. \textbf{The} tar is from \textbf{the} roof."So they really were working under tough conditions.
Here,"A mouse has climbed into \textbf{the} blower behind \textbf{the} regulator rack, set blower to vibrating.
Result : no more mouse.""Here lies mouse.
Born :?.
Died : 4:50 a. m., May 1953."There ' s an inside joke someone has penciled in :"Here lies Marston Mouse."If you ' re a mathematician, you get that, because Marston was a mathematician who \textbf{objected} to \textbf{the} \textbf{computer} being there."Picked a lightning bug off \textbf{the} \textbf{drum}.""Running at two kilocycles."That ' s two thousand cycles per second —"yes, I ' m chicken"— so two kilocycles was slow speed. \textbf{The} high speed was 16 kilocycles.
I don ' t know if you remember a Mac that was 16 Megahertz, that ' s slow speed."I have now duplicated both results.
How will I know which is right, assuming one result is correct?
This now is \textbf{the} third different output.
I know when I ' m licked.""We ' ve duplicated errors before.""Machine run, fine.
Code isn ' t.""Only happens when \textbf{the} machine is running."And sometimes things are \textbf{okay}."Machine a thing of beauty, and a joy forever.""Perfect running.""Parting thought : when there ' s bigger and better errors, we ' ll have them."So, nobody was supposed to know they were actually designing bombs.
They ' re designing hydrogen bombs.
But someone in \textbf{the} logbook, late one night, finally drew a bomb.
So, that was \textbf{the} result.
It was Mike, \textbf{the} first thermonuclear bomb, in 1952.
That was designed on that machine, in \textbf{the} woods behind \textbf{the} Institute.
So Von Neumann invited a whole gang of weirdos from all over \textbf{the} world to work on all these problems.
Barricelli, he came to do what we now call, really, \textbf{artificial} life, trying to see if, in this \textbf{artificial} \textbf{universe} — he was a viral-geneticist, way, way, way ahead of his time.
He ' s still ahead of some of \textbf{the} stuff that ' s being done now.
Trying to start an \textbf{artificial} genetic system running in \textbf{the} \textbf{computer}.
Began — his \textbf{universe} started March 3, ' 53.
So it ' s almost exactly — it ' s 50 years ago next Tuesday, I guess.
And he saw everything in terms of — he could read \textbf{the} binary code straight off \textbf{the} machine.
He had a wonderful rapport.
Other people couldn ' t get \textbf{the} machine running.
It always worked for him.
Even errors were duplicated."Dr.
Barricelli claims machine is wrong, code is right."So he designed this \textbf{universe}, and ran it.
When \textbf{the} bomb people went home, he was allowed in there.
He would run that thing all night long, running these things, if anybody remembers Stephen Wolfram, who reinvented this stuff.
And he published it.
It wasn ' t locked up and \textbf{disappeared}.
It was published in \textbf{the} literature."If it ' s that easy to create living organisms, why not create a few yourself?"So, he decided to give it a try, to start this \textbf{artificial} \textbf{biology} going in \textbf{the} machines.
And he found all these, sort of — it was like a naturalist coming in and looking at this tiny, 5, 000-byte \textbf{universe}, and seeing all these things happening that we see in \textbf{the} outside world, in \textbf{biology}.
This is some of \textbf{the} generations of his \textbf{universe}.
But they ' re just going to stay numbers ; they ' re not going to become organisms.
They have to have something.
You have a genotype and you have to have a phenotype.
They have to go out and do something.
And he started doing that, started giving these little numerical organisms things they could play with — playing chess with other machines and so on.
And they did start to \textbf{evolve}.
And he went around \textbf{the} country after that.
Every time there was a new, fast machine, he started using it, and saw exactly what ' s happening now.
That \textbf{the} programs, instead of being turned off — when you quit \textbf{the} program, you ' d keep running and, basically, all \textbf{the} sorts of things like Windows is doing, running as a multi-cellular organism on many machines, he envisioned all that happening.
And he saw that \textbf{evolution} itself was an intelligent process.
It wasn ' t any sort of creator intelligence, but \textbf{the} thing itself was a giant parallel computation that would have some intelligence.
And he went out of his way to say that he was not saying this was lifelike, or a new kind of life.
It just was another version of \textbf{the} same thing happening.
And there ' s really no difference between what he was doing in \textbf{the} \textbf{computer} and what nature did billions of years ago.
And could you do it again now?
So, when I went into these archives looking at this stuff, lo and behold, \textbf{the} archivist came up one day, saying,"I think we found another box that had been thrown out."And it was his \textbf{universe} on punch cards.
So there it is, 50 years later, sitting there — sort of suspended animation.
That ' s \textbf{the} instructions for running — this is actually \textbf{the} source code for one of those \textbf{universes}, with a note from \textbf{the} engineers saying they ' re having some problems."There must be something about this code that you haven ' t explained yet."And I think that ' s really \textbf{the} truth.
We still don ' t understand how these very simple instructions can lead to increasing complexity.
What ' s \textbf{the} dividing line between when that is lifelike and when it really is alive?
These cards, now, thanks to me showing up, are being saved.
And \textbf{the} question is, should we run them or not?
You know, could we get them running?
Do you want to let it loose on \textbf{the} Internet?
These machines would think they — these organisms, if they came back to life now — whether they ' ve died and gone to heaven, there ' s a \textbf{universe}.
My \textbf{laptop} is 10 thousand million times \textbf{the} size of \textbf{the} \textbf{universe} that they lived in when Barricelli quit \textbf{the} project.
He was thinking far ahead, to how this would really grow into a new kind of life.
And that ' s what ' s happening!
When Juan Enriquez told us about these 12 trillion bits being transferred back and forth, of all this genomics data going to \textbf{the} proteomics lab, that ' s what Barricelli imagined : that this digital code in these machines is actually starting to code — it already is coding from nucleic acids.
We ' ve been doing that since, you know, since we started PCR and synthesizing small strings of DNA.
And real soon, we ' re actually going to be synthesizing \textbf{the} proteins, and, like Steve showed us, that just opens an entirely new world.
It ' s a world that Von Neumann himself envisioned.
This was published after he died : his sort of unfinished notes on self-reproducing machines, what it takes to get \textbf{the} machines sort of jump-started to where they begin to reproduce.
It took really three people : Barricelli had \textbf{the} concept of \textbf{the} code as a living thing ; Von Neumann saw how you could build \textbf{the} machines — that now, last count, four million of these Von Neumann machines is built every 24 hours ; and Julian Bigelow, who died 10 days ago — this is John Markoff ' s obituary for him — he was \textbf{the} important missing link, \textbf{the} engineer who came in and knew how to put those vacuum tubes together and make it work.
And all our \textbf{computers} have, inside them, \textbf{the} \textbf{copies} of \textbf{the} architecture that he had to just design one day, sort of on pencil and paper.
And we owe a tremendous credit to that.
And he explained, in a very generous way, \textbf{the} spirit that brought all these different people to \textbf{the} Institute for Advanced Study in \textbf{the} ' 40s to do this project, and make it freely available with no \textbf{patents}, no restrictions, no intellectual property disputes to \textbf{the} rest of \textbf{the} world.
That ' s \textbf{the} last entry in \textbf{the} logbook when \textbf{the} machine was shut down, July 1958.
And it ' s Julian Bigelow who was running it until midnight when \textbf{the} machine was officially turned off.
And that ' s \textbf{the} end.
Thank you very much.

% matched lemmas: abstract, array, artificial, biology, clon, computer, copy, disappear, disk, drink, drum, electrical, evolution, evolve, geek, graphics, laptop, object, oil, okay, patent, slide, technique, the, universe
\end{textsample}
